\documentclass[a4paper,11pt]{article}

\usepackage[T1]{fontenc}
\usepackage[utf8]{inputenc}
\usepackage[ngerman]{babel}
\usepackage{geometry}
\geometry{margin=2.2cm}
\usepackage{booktabs}
\usepackage{longtable}
\usepackage{xcolor}
\usepackage{listings}
\usepackage{tikz}
\usetikzlibrary{arrows.meta,positioning,shapes.geometric}
\usepackage{hyperref}
\usepackage{float}

\lstdefinestyle{shared}{
  basicstyle=\ttfamily\small,
  keywordstyle=\color{blue!60!black},
  commentstyle=\color{green!40!black},
  stringstyle=\color{orange!60!black},
  numbers=left,
  numberstyle=\scriptsize\color{gray},
  stepnumber=1,
  numbersep=8pt,
  showstringspaces=false,
  breaklines=true,
  frame=single
}

\title{Technische Dokumentation der Dashboards\\
\large Datenvisualisierungsprojekt (WS 2022/2023)}
\author{AFL et al.}
\date{\today}

\begin{document}
\maketitle
\tableofcontents
\newpage

\section{Zielsetzung}
Dieses Projekt vereint zwei interaktive Analyseanwendungen:
\begin{itemize}
  \item ein Python/Dash-Dashboard für Einzelhandelsumsätze (\texttt{Dashboard-01}),
  \item ein R/Shiny-Dashboard für den europäischen Flugverkehr (\texttt{Dashboard-02}).
\end{itemize}

Die überarbeitete Version fokussiert auf Robustheit, reproduzierbare Ausführung und bessere Wartbarkeit.

\section{Dateiübersicht und Aufgaben}
\begin{longtable}{p{0.35\linewidth} p{0.58\linewidth}}
\toprule
\textbf{Datei} & \textbf{Funktion} \\
\midrule
\texttt{Dashboard-01/app.py} & Python-Dash-Anwendung mit Monatsvergleich, KPI-Karten und Store-Ranking. \\
\texttt{Dashboard-01/*.csv} & Datengrundlage für Monats-, Wochen-, Feiertags- und Filialumsätze. \\
\texttt{Dashboard-02/db.R} & R-Shiny-Anwendung mit reaktiver Länderfilterung des Flugverkehrs. \\
\texttt{Dashboard-02/fluege.csv} & Zeitreihe der wöchentlichen Fluganzahl pro Land. \\
\texttt{README.md} & Projektdokumentation mit Setup- und Startanleitung. \\
\bottomrule
\end{longtable}

\section{Dashboard-01: Python/Dash}
\subsection{Funktionale Beschreibung}
Die Anwendung lädt mehrere Umsatzdatensätze, prüft die benötigten Spalten und erzeugt daraus vier UI-Bereiche:
\begin{itemize}
  \item Monat-Auswahl (Basis und Referenz),
  \item KPI-Karte "Gesamtumsatz",
  \item KPI-Karte "Ferienumsatz",
  \item Vergleichsvisualisierungen für Wochenverlauf und Top-Filialen.
\end{itemize}

\subsection{Verbesserungen}
\begin{itemize}
  \item Moderne Dash-Imports mit Fallback für ältere Dash-Versionen.
  \item Dateipfade relativ zur Skriptdatei statt abhängig vom aktuellen Arbeitsverzeichnis.
  \item Validierung der notwendigen Spalten pro Datensatz.
  \item Sichere Monatsauswahl und robuste Berechnung der Kennzahlen (ohne fragile Indexzugriffe).
  \item Bereinigte Plot-Erstellung mit klaren Hilfsfunktionen.
\end{itemize}

\subsection{Ablaufdiagramm (Python)}
\begin{figure}[H]
\centering
\begin{tikzpicture}[
  node distance=1.2cm and 1.8cm,
  block/.style={rectangle, draw, rounded corners, align=center, minimum width=4.8cm, minimum height=0.95cm},
  line/.style={-Latex, thick}
]
\node[block] (csv) {CSV-Dateien laden\\(\texttt{monthly}, \texttt{weekly}, \texttt{store}, \texttt{holiday}, \texttt{dept})};
\node[block, below=of csv] (validate) {Spalten prüfen und Datentypen harmonisieren};
\node[block, below=of validate] (ui) {Layout initialisieren\\(Navbar, Dropdowns, Karten)};
\node[block, below=of ui] (callback) {Callback: \texttt{update\_cards(base, comparison)}};
\node[block, below left=of callback, xshift=1.2cm] (metrics) {KPI-Berechnung\\Gesamt- und Ferienumsatz};
\node[block, below right=of callback, xshift=-1.2cm] (plots) {Plot-Berechnung\\Wochenvergleich + Top Stores};
\node[block, below=of callback, yshift=-2.2cm] (render) {Rückgabe der vier Karten an das Frontend};

\draw[line] (csv) -- (validate);
\draw[line] (validate) -- (ui);
\draw[line] (ui) -- (callback);
\draw[line] (callback) -- (metrics);
\draw[line] (callback) -- (plots);
\draw[line] (metrics) -- (render);
\draw[line] (plots) -- (render);
\end{tikzpicture}
\caption{Daten- und Kontrollfluss im Dash-Callback.}
\end{figure}

\subsection{Visualisierung der Monatsdaten}
\begin{figure}[H]
\centering
\begin{tikzpicture}[x=1cm,y=0.04cm]
\draw[->,thick] (0,0) -- (13.6,0) node[right] {Monat};
\draw[->,thick] (0,0) -- (0,190) node[above] {Umsatz [Mio.]};

\foreach \m/\v [count=\i from 1] in {
  Jan/78.2,Feb/94.5,Mar/86.2,Apr/108.3,May/84.4,Jun/92.1,
  Jul/108.6,Aug/94.6,Sep/105.1,Oct/88.9,Nov/113.8,Dec/165.3
} {
  \filldraw[fill=blue!45,draw=blue!70!black] (\i-0.33,0) rectangle (\i+0.33,\v);
  \node[below,rotate=45,anchor=east] at (\i,0) {\scriptsize \m};
}

\draw[red,thick] plot[smooth] coordinates {
  (1,0.0) (2,24.1) (3,0.0) (4,0.0) (5,0.0) (6,0.0)
  (7,0.0) (8,0.0) (9,23.0) (10,0.0) (11,41.2) (12,20.1)
};
\foreach \x/\y in {
  1/0.0,2/24.1,3/0.0,4/0.0,5/0.0,6/0.0,
  7/0.0,8/0.0,9/23.0,10/0.0,11/41.2,12/20.1
} {
  \fill[red] (\x,\y) circle (1.6pt);
}

\node[blue!70!black] at (10.2,176) {\scriptsize Balken: Gesamtumsatz};
\node[red!80!black] at (10.2,166) {\scriptsize Linie: Ferienumsatz};
\end{tikzpicture}
\caption{Monatsumsatz und Ferienumsatz aus \texttt{monthly\_sales\_df.csv}.}
\end{figure}

\subsection{Codeauszug (Python)}
\begin{lstlisting}[style=shared,language=Python]
from pathlib import Path
import pandas as pd

BASE_DIR = Path(__file__).resolve().parent

def load_dataframe(file_name: str) -> pd.DataFrame:
    path = BASE_DIR / file_name
    df = pd.read_csv(path)
    unnamed = [c for c in df.columns if str(c).startswith("Unnamed")]
    if unnamed:
        df = df.drop(columns=unnamed)
    return df

def month_metric(month_name: str, column_name: str) -> float:
    values = monthly_sales_df.loc[monthly_sales_df["Month"] == month_name, column_name]
    return float(values.iloc[0]) if not values.empty else 0.0
\end{lstlisting}

\section{Dashboard-02: R/Shiny}
\subsection{Funktionale Beschreibung}
Die Shiny-Anwendung lädt Flugdaten, konvertiert Zeitstempel, erlaubt eine Länderselektion und zeigt die wöchentliche Gesamtzahl als Balkendiagramm.

\subsection{Verbesserungen}
\begin{itemize}
  \item Robuste Pfadauflösung über Arbeitsverzeichnis und Skriptpfad.
  \item Konsistente Datumskonvertierung (\texttt{\%Y-\%m-\%dT\%H:\%M:\%SZ}).
  \item Reaktive Filterung mit Validierung (Hinweis bei leerer Auswahl).
  \item Verbessertes Plot-Layout (Achsenformatierung, lesbare Datumsbeschriftung, Legende unten).
\end{itemize}

\subsection{Ablaufdiagramm (R/Shiny)}
\begin{figure}[H]
\centering
\begin{tikzpicture}[
  node distance=1.15cm and 1.8cm,
  block/.style={rectangle, draw, rounded corners, align=center, minimum width=4.8cm, minimum height=0.95cm},
  line/.style={-Latex, thick}
]
\node[block] (read) {CSV laden: \texttt{fluege.csv}};
\node[block, below=of read] (clean) {Datum und numerische Werte bereinigen};
\node[block, below=of clean] (ui) {UI aufbauen\\(Titel, Checkboxen, Plotbereich)};
\node[block, below=of ui] (reactive) {Reaktiv: Länderfilter\\\texttt{plot\_df()}};
\node[block, below=of reactive] (render) {\texttt{renderPlot()} mit \texttt{ggplot2}};

\draw[line] (read) -- (clean);
\draw[line] (clean) -- (ui);
\draw[line] (ui) -- (reactive);
\draw[line] (reactive) -- (render);
\end{tikzpicture}
\caption{Reaktiver Ablauf in der Shiny-App.}
\end{figure}

\subsection{Visualisierung: Durchschnittliche Flüge je Land}
Die folgende Grafik zeigt die durchschnittliche wöchentliche Fluganzahl je Land (aus \texttt{fluege.csv}, über den gesamten Zeitraum).

\begin{figure}[H]
\centering
\begin{tikzpicture}[x=1.6cm,y=0.0002cm]
\draw[->,thick] (0,0) -- (7.2,0) node[right] {Land};
\draw[->,thick] (0,0) -- (0,36000) node[above] {Ø Flüge/Woche};

\foreach \l/\v [count=\i from 1] in {
  Belgium/5425.73,
  Frankreich/29704.58,
  Ireland/4403.48,
  Luxemburg/1194.66,
  Niederlande/8804.97,
  UK/31463.48
} {
  \filldraw[fill=teal!45,draw=teal!70!black] (\i-0.3,0) rectangle (\i+0.3,\v);
  \node[below,rotate=20,anchor=east] at (\i,0) {\scriptsize \l};
}
\end{tikzpicture}
\caption{Aggregierte Kennzahl aus den Flugdaten (über den kompletten Datensatz).}
\end{figure}

\subsection{Codeauszug (R)}
\begin{lstlisting}[style=shared,language=R]
flights <- read_csv(resolve_data_path("fluege.csv"), show_col_types = FALSE) %>%
  mutate(
    Date = as.Date(Date, format = "%Y-%m-%dT%H:%M:%SZ"),
    Total = as.numeric(Total)
  ) %>%
  filter(!is.na(Date), !is.na(Total)) %>%
  arrange(Date)

plot_df <- reactive({
  req(input$land)
  flights %>% filter(Land %in% input$land)
})
\end{lstlisting}

\section{Zusammenfassung}
Die überarbeitete Version erhöht Stabilität und Lesbarkeit beider Anwendungen deutlich:
\begin{itemize}
  \item reproduzierbare Dateipfade und robuste Dateneinleseprozesse,
  \item sauberere und defensivere Callback-/Reactive-Logik,
  \item nachvollziehbare, strukturierte Visualisierungslogik,
  \item klare Dokumentation für Betrieb und Weiterentwicklung.
\end{itemize}

\end{document}
